\documentclass{article}
\usepackage[margin=2cm]{geometry}
\usepackage{amsmath}
\usepackage{graphicx}
\usepackage{booktabs}
\usepackage[utf8]{inputenc}
\usepackage{ragged2e}
\setlength{\emergencystretch}{3em}
\begin{document}
\sloppy
\begin{itemize}
\item 10. Proof. We omit (a) since is standard. For (b), if u attains an interior maximum, then the
conclusion follows from strong maximum principle.

\end{itemize}

If not then for some xE aUux)> u(x) Vx E U. Then Hopf's lemma implies x)>0,
which is a contradiction.


Remark 2. A generalization of this problem to mixed boundary conditions is recorded in
Gilbarg-Trudinger, Elliptic PDEs of second order, Problem 3.1.


11.Proof.Define


\begin{center}
\includegraphics[width=0.5\linewidth]{formula_001.png}
\end{center}


By Exercise 5.17uE HU.Thenfor all v E CUu0


\begin{equation}
\begin{array} { r l } { B [ \hat { \phi } ( u ) , v ] = \int _ { U } \sum _ { i , j } a ^ { i j } ( \hat { \phi } ( u ) ) _ { x _ { i } } v _ { x _ { j } } ~ d x } & { } \\ { = } & { \int _ { U } \sum _ { i , j } a ^ { i j } \hat { \phi } ^ { \prime } ( u ) u _ { x _ { i } } v _ { x _ { j } } ~ d x , ~ ( \phi ^ { \prime } ( u ) \mathrm { ~ i s ~ b o u n d e d ~ s i n c e ~ } u \mathrm { ~ i s ~ b o u n d e d } ) } \\ { } & { = \int _ { U } \sum _ { i , j } a ^ { i j } u _ { x _ { i } } ( \hat { \phi } ^ { \prime } ( u ) v ) _ { x _ { i } } - \sum _ { i , j } a _ { i j } \hat { \phi } ^ { \prime \prime } ( u )}\end{array}
\end{equation}


(We don't know whether the product of two H' functions is weakly differentiable. This is why.
we do not take v E H.) Now we complete the proof with the standard density argument.


12. Proof. Given u E C2(U) n C(U) with Lu  0 in U and u  0 on aU. Since U is compact and
ECU)c>0.So w=EC2UnCU.Brutal computation gives us


\begin{equation}
\begin{array} { r l } { - a ^ { i j } w _ { z _ { i } x _ { j } } } & { = \frac { - a ^ { i j } u _ { z _ { i } x _ { j } } v + a ^ { i j } v _ { z _ { i } x _ { j } } u } { u ^ { 2 } } + \frac { a ^ { i j } v _ { z _ { i } } u _ { z _ { j } } - a ^ { i j } u _ { z _ { i } } v _ { z _ { j } } } { v ^ { 2 } } - a ^ { i j } \frac { 2 } { v } v _ { z _ { j } } \frac { u _ { z _ { i } } u - v u _ { z _ { i } } } { v ^ { 2 } } } \\ & { \qquad }\end{array}
\end{equation}


Therefore,


\begin{equation}
\begin{array} { r l r } & { } & { M w : = - a ^ { i j } w _ { z _ { i } x _ { j } } + w _ { z _ { i } } \big [ b ^ { i } - a ^ { i j } \frac { 2 } { v } v _ { z _ { j } } \big ] = \frac { L u } { v } - \frac { u L v } { v ^ { 2 } } \leq 0 ~ \mathrm { ~ o n ~ } \{ x \in \bar { U } : u > 0 \} \subseteq U } \\ & { } & { \qquad \dots \quad \dots \quad } \\ & { } & { \qquad \dots \quad \dots \quad } \end{array}
\end{equation}


If \{x E U : u > 0\} is not empty, Weak maximum principle to the operator M with bounded
coefficeints (since v E C (U)) will lead a contradiction that


\begin{equation}
0 < \operatorname* { m a x } _ { \{ u > 0 \} } w = \operatorname* { m a x } _ { \partial \{ u > 0 \} } w = \frac { 0 } { v } = 0
\end{equation}


Hence u  0 in U.

\end{document}
